\documentclass[11pt]{article}
\usepackage[utf8]{inputenc}
\usepackage{amsmath, amssymb, amsthm}
\usepackage{graphicx}
\usepackage{geometry}
\usepackage{hyperref}
\usepackage{tabularx}
\usepackage{booktabs}
\geometry{a4paper, margin=1in}
%\usepackage{tikz}
%\usetikzlibrary{calc, positioning, arrows.meta}
\usepackage{listings}
\usepackage{graphicx}
\usepackage{xcolor}


\definecolor{codegray}{rgb}{0.5,0.5,0.5}
\definecolor{codepurple}{rgb}{0.58,0,0.82}
\definecolor{backcolour}{rgb}{0.95,0.95,0.92}
\lstset{
    language=Python,
    backgroundcolor=\color{backcolour},
    commentstyle=\color{codegray},
    keywordstyle=\color{magenta},
    stringstyle=\color{codepurple},
    basicstyle=\ttfamily\small,
    breaklines=true,
    showstringspaces=false
}

\title{Programming Lab: Optimizing Python with GPU}
\author{VB}
\date{\today}

\begin{document}

\maketitle



\section*{What to install}


You will need Python 3.8+ and the following standard libraries.
Open your terminal and run:
\begin{lstlisting}
    pip install numpy matplotlib pyvista vispy PyQt5
\end{lstlisting}

To actually write the fast GPU code, you need CuPy. CuPy requires an Nvidia GPU and the CUDA toolkit installed on your system.
\begin{lstlisting}
    pip install cupy
\end{lstlisting}


\subsection*{\textcolor{red}{Read this if you cannot run Nvidia!}}

If you are on a Mac or a computer without a dedicated Nvidia GPU, \lstinline|import cupy| will crash. To complete the assignments, simply change the import statement at the top of each file to: \lstinline|import numpy as cp|.

Your code will run, you will learn the vectorized math, and you \textbf{will be graded normally (as if it was cupy)}, but you will not get the massive speed boost when you turn off debug mode.


\section*{Problem organization}

You will get a number of numpy files with \textbf{already working code}\,--- use this fact to your advantage, e.g. compare if your code produces the same result during debug.
\begin{itemize}
 \item \lstinline|game_of_life.py| (Game of Life, 2D Array Slicing and Boolean Logic).
\item \lstinline|galaxies_colliding.py| (Galaxy Collision, N-Body problems, $O(N^2)$ bottlenecks, and Matrix Broadcasting).
\item \lstinline|lennard_jones.py| (Thermodynamics, Complex math kernels, Thermostats, and Boundary Collisions).
\item \lstinline|mandelbrot.py| (Fractals, In-place data generation and Escape-time Boolean masking).
\item \lstinline|waves_3d.py| (3D Wave Equation, Partial Differential Equations and Time-History Buffers).
\end{itemize}
Inside of each file you will find functions designated as subproblems. You have to look for blocks like these:
\begin{lstlisting}
# ===================================================================
# SUBPROBLEM 1: COMPUTE LENNARD-JONES FORCES (The O(N^2) Bottleneck)
# Optimize this function removing cycles
# ===================================================================
\end{lstlisting}
Your goal is to optimize these functions removig Python cycles and using CuPy/numpy arrays. For more details see comments in appropriate files.

\textbf{Note:} all given files can be run immediately and you can see some result, but they use slow-machine-parameters (e.g. low number particles). If you successfully optimize the code\,--- change \lstinline|DEBUG = True| to \lstinline|DEBUG = False| and you should see nice demo with large number of particles.


\section*{\textcolor{red}{Submission}}

Submit all files listed above for review.


\end{document}

